% Source-derived chapter 03: Contractions of surfaces
% Original source: no-enriques-surface-over-the-integers
\chapter{Contractions of surfaces}
\label{no-enriques-surface-over-the-integers-chapter-03}
\source{no-enriques-surface-over-the-integers}

\begin{remark}[Source section]
\label{no-enriques-surface-over-the-integers-chapter-03-source}
\source{no-enriques-surface-over-the-integers}
\source{no-enriques-surface-over-the-integers}
This chapter reproduces the corresponding section of the retrieved
LaTeX source, including its definitions, statements, examples, and
proofs. It has not yet been translated into Lean.
\notready
\end{remark}

% The source section title was: Contractions of surfaces
\mylabel{Contractions}

We now examine in dimension two the results of the preceding section.
Let $k$ be a ground field of characteristic $p\geq 0$, and 
 $S$ be a smooth proper surface with $h^0(\O_S)=1$.
Let $C\subset X$ be a curve. Decompose it into irreducible components  $C=m_1C_1+\ldots+m_rC_r$ with multiplicities $m_i\geq 1$, 
and let 
$$
N=N(C)=(C_i\cdot C_j)_{1\leq i,j\leq r}\in\Mat_r(\ZZ)
$$
be the resulting intersection matrix.  
We say that the curve $C$ is \emph{negative-definite}  or \emph{negative-semidefinite} if the
intersection matrix $N$ has the respective property. 
One may use the intersection numbers  to endow $\Gamma=\Gamma(C)$ with the
structure of an \emph{edge-labeled graph}, where the   labels attached to the edges are the intersection numbers
$(C_i\cdot C_j)>0$ for $i\neq j$

Now let  $f:S\ra Z$ be a
contraction. Then  either $Z$ is a geometrically normal   surface and the morphism is birational,
or $Z$ is a smooth curve and the morphism is of fiber type.
Let $\ba:\Spec(\Omega)\ra Z$ be a geometric point, for some algebraically closed field $\Omega$,
with image point $a\in Z$.
Write $S_\ba=X\times_Z\Spec(\Omega)=f^{-1}(\ba)$ and $S_a=S\times_Z\Spec\kappa(a)=f^{-1}(a)$ for the resulting geometric 
and schematic fibers, respectively. 
The goal of this section is to establish the following:

\begin{theorem}
\source{no-enriques-surface-over-the-integers}
\notready
\mylabel{geometric and schematic fiber}
Suppose the geometric fiber $S_\ba$   is one-dimensional.
If the group scheme $\Num_{S/k}$ is constant,   the following holds:
\begin{enumerate}
\item 
The map $\Gamma(S_\ba)\ra\Gamma(S_a)$ is a graph isomorphism.
\item 
The above graph isomorphism respects edge labels, provided that the  reduced scheme  $(S_a)_\red$ 
is geometrically reduced  over the residue field $\kappa(a)$.
\item 
If the geometric fiber $S_\ba$ is reducible, the residue field extension $k\subset\kappa(a)$ is purely inseparable.
\end{enumerate}
\end{theorem}


\proof
Assertion (ii) is an immediate consequence from (i) and the definitions, and left to the reader.
To verify (i) and (iii) we write  $E=\kappa(a)$. Without restriction, it suffices to treat the case that $\Omega= E^\alg$.
The geometric fiber $S_\ba=S_a\otimes_E\Omega\subset S_a\otimes_k\Omega$ is a   curve on the base-change $S\otimes_k\Omega$.
This curve is connected, by the condition $\O_Z=f_*(\O_S)$.
Let  $C_i\subset S_\ba$, $1\leq i\leq r$ be the irreducible components, and write $N=(C_i\cdot C_j)_{1\leq i,j\leq r}$
for the resulting intersection matrix.

We first consider   the case that $f:X\ra Z$ is birational. Then $N$ is negative-definite, in particular
the canonical map $\bigoplus\ZZ C_i\ra \Num_{S/k}(\Omega)$ is injective. It follows that the Galois action
on the set $\{C_1,\ldots,C_r\}$ is trivial. In turn, $\Gamma(S_\ba)\ra\Gamma(S_a)$ is a graph isomorphism.
Now suppose that $r\geq 2$. Seeking a contradiction, we assume that $k\subset E$ is not purely inseparable.
Then $S_a\otimes_k\Omega$ is disconnected. On the other hand, the fiber $S_a$ is connected,
 by the condition $\O_Z=f_*(\O_S)$.
It follows that the Galois action on the   irreducible components $E_i\subset S_a\otimes_k\Omega$
is non-trivial.
As above, the canonical map $\bigoplus\ZZ E_i\ra\Num_{S/k}(\Omega)$ is injective, so the Galois action on 
the set $\{E_1,\ldots,E_s\}$ must be trivial, contradiction.

It remains to treat the case that  $f:X\ra Z $ is of fiber type. Then the matrix $N$ is negative-semidefinite.
For $r=1$ our graphs have but one vertex, and the map $\Gamma(S_\ba)\ra\Gamma(S_a)$ is obviously a graph isomorphism.
Now suppose  $r\geq 2$. Then  the curves $C_1+\ldots+C_{r-1}$ and $C_2+\ldots+C_r$ are negative-definite. Using the 
previous paragraph,
we infer that $\Gamma(S_\ba)\ra\Gamma(S_a)$ is again a graph isomorphism. Likewise, one shows that
 $k\subset E$ is purely inseparable.
\qed

 
%===========================================================
